\section{Summarizing Data Structure} 
\label{sec:summarizing_data_structure}
Generally there are two main catalogues for summarizing data structure: the hierarchical catalogue and the partitional catalogue. In the following, we discuss six representative summarizing data structures commonly used in \dsc algorithms. 
Three are hierarchical-based, and three are partitional-based. Their main features is summarized in Figure~\ref{fig:data_structure}.

\begin{figure}[t] 
\centering
\includegraphics[width=0.7\textwidth]{Figures/Design Aspects/Summarizing Data Structure.pdf}
\caption{Six Types of Summarizing Data Structure.}
\label{fig:data_structure}
\end{figure}

\subsection{Hierarchical Summarizing Data Structure}
\label{subsec:hierarchical}
Hierarchical summarizing data structure groups the incoming data points into a tree structure. As shown in Figure~\ref{fig:data_structure}(a), each temporal cluster is depicted as a \emph{node}, and the parent cluster comprises multiple sub-clusters within the tree.
% Every parent node $P$ in this type of structure represents a bigger cluster and can be divided into some small sub-clusters represented by the child nodes of $P$. Therefore,
To insert a new data point, it starts by searching for a suitable temporal cluster from the root node to one of the leaf nodes. Such a tree structure makes the insertion fast as there is no need for checking every existing temporal cluster. In the following, we discuss three popular and representative hierarchical summarizing data structures different from each other for the implementation of each tree node.
% , including clustering feature tree (\CFT), coreset tree (\CoreT), and dependency tree (\DPT).


\textbf{Clustering Feature Tree (\CFT).}
\label{subsec:cf-tree}
\CFT~\cite{BIRCH:96} is one of the oldest but representative hierarchical summarizing data structures. Every tree node of \CFT is named \emph{clustering feature (CF)}, which consists of
% only consists of three elements recording 
the position information of a temporal cluster. Although CF is quite simple, it is adequate to support many basic operations of stream clustering, such as distance calculation and cluster updating. Due to the additivity of the CF~\cite{BIRCH:96}, \CFT can be updated incrementally rather than starting from scratch again. Moreover, it also offers some advanced operations such as merging or splitting for handling \cef during stream clustering.

\textbf{Coreset Tree (\CoreT).}
\label{subsec:CT}
\CoreT~\cite{StreamKM++:12} is a binary tree structure used for extracting a core subset from a large volume of stream data. Compared with CF in \CFT, the tree node of \CoreT named \emph{CoreNode} provides additional information about the density of the corresponding cluster. However, since \CoreT is unable to update incrementally like \CFT, it needs to rebuild the whole tree structure during the clustering procedure, which might influence the clustering efficiency of the whole algorithm. To overcome this shortage, \CoreT reduces the time of rebuilding via lazily performing the update. In exchange, it may not be able to provide real-time clustering information for accurate clustering, especially when \cef happens frequently.

\textbf{Dependency Tree (\DPT).}
\label{subsec:DPT}
\DPT~\cite{EDMStream:17} is one of the most recent hierarchical summarizing data structures. It is specifically designed for handling \cef.
% all types of \cef 
% as discussed in section~\ref{sec:background}. 
The tree node of \DPT is called \emph{cluster cell} (CC).
\DPT is built based on the cluster density and it requires every cluster with low density to be connected to its nearest denser cluster. Although CC contains similar clustering information as CoreNode in \CoreT, it can partially adjust itself for adapting to the evolving activities in the stream, helping to achieve a higher clustering efficiency.
% Although CC is not as flexible as CF, it additionally provides the temporal density of every cluster, helping to achieve a higher clustering accuracy, .
% can not only be used to locate the corresponding cluster but also provides the temporal density of every cluster for more precise clustering. This may lead to a higher clustering accuracy.}
% due to using more comprehensive information

\subcompact
\subsection{Partitional Summarizing Data Structure}
\label{subsec:partition}
As shown in Figure~\ref{fig:data_structure}(b), partitional
% the other catalogue of 
summarizing data structures\margi{R1O1}do not organize temporal clusters into hierarchical structures. They hence denote temporal clusters as \emph{objects}, rather than nodes.
% is based on partition, which 
% treats every cluster as an individual object instead of integrating them into a tree like the hierarchical-based one. 
When inserting a new data point, they need to check every object. This brings significant overhead compared to hierarchical structures but may help to achieve a higher clustering accuracy.
% Therefore, it will not be as efficient as the hierarchical-based one but may bring in a better clustering accuracy. 
In the following, we discuss three typical partitional summarizing data structures, different from each other by the implementation of objects.
% including micro clusters (\MCs), grids (\Grids), and meyerson sketch (\MS).


% Therefore, every time in order to find the belonging cluster for inserting the new data point, every existing cluster may have to be checked, which will bring in a more precise cluster but may also lead to a poor clustering efficiency. }
% The use of \emph{clustering feature tree} (\CFT) for summarizing large amounts of data was first introduced in the \birch algorithm. The \CFT is a very appealing summary structure for streams as it can be easily maintained online, thanks to its incremental and additive properties~\cite{Survey13}.

% As shown in Figure~\ref{fig:data_structure}(a), the \CFT is a height-balanced tree, where each non-leaf node contains at most $B$ entries, and each leaf node contains at most $L$ entries, where each entry is a CF.
% % , having each a CF and a pointer to a child node. Every leaf node contains at most $L$ entries, where each entry is a CF. 
% Each CF=($N$, $LS$, $SS$) has three components: 
% $N$, the number of data objects, 
% $LS$, the linear sum of the data objects, 
% and $SS$, the sum of squared data objects. 
% The structures $LS$ and $SS$ are n-dimensional arrays. These three components are used to compute cluster measures, such as cluster mean ($\frac{LS}{N}$), radius ($\sqrt{(\frac{SS}{N} - (\frac{LS}{N})^2)}$), and diameter ($\sqrt{(\frac{2N*SS-2*LS^2}{N(N-1)})}$).
% Each leaf node has a maximum diameter (or radius) represented by a user-defined parameter, $T$. The value of this parameter defines whether a new data object may be absorbed by a CF. Thus, $T$ determines the size of the CF-Tree, where higher values of $T$ lead to smaller trees.
 
% When a new data point arrives, it descends the \CFT from the root to the leaves by choosing in each non-leaf node its closest CF entry (closeness is defined by the Euclidean distance between new objects and the centroids of CF in non-leaf nodes). In a leaf node, the closest entry is selected and tested to verify whether it can absorb the new object. If so, the CF is updated, otherwise, a new CF is created containing this particular data point. 

\textbf{Micro Clusters (\MCs).} 
\label{subsec:MCs}
\MCs~\cite{Clustream:03} is one of the oldest partitional summarizing data structures. It contains the same positional information as CF in \CFT. To keep track of the clusters' activities in real-time, it has two additional elements used for summarizing the timestamps for the clusters' updates. Therefore, it is more comprehensive and robust for stream clustering with high accuracy, especially under \cef.
% some key characteristics such as \cef compared with \CFT.

\textbf{Grids (\Grids).} 
\label{subsec:Grids}
\Grids~\cite{DStream:2007} performs its clustering simply based on checking whether the data lies in a specific grid of $d$-dimensional space and grouping the dense grids into temporal clusters. 
% During new point insertions, density, time, and sparsity information of every grid needs to be updated. In exchange, 
Compared to other data structures, \Grids is more computationally efficient because it does not require frequent distance calculations between new data and the grid. Additionally, it periodically removes sparse grids, which further boosts its clustering efficiency. However, since every grid has a fixed position in the space and the information about the relationship among grids is only updated lazily, \Grids may not work well when some evolving activities occur frequently among the grid and thus leads to lower clustering accuracy.

\textbf{Augmented Meyerson Sketch (\MS).} 
\label{subsec:meyerson_sketch}
\MS~\cite{SL-KMeans:20} is the extension of traditional Meyerson Sketch~\cite{MeyersonSketch:01} which can better summarize the clustering information using the limited number of data points. \MS mainly consists of a set of temporally weighted centres and the weight is mainly influenced by both the time and density information of the corresponding cluster. 
Similar to \CoreT, \MS may need to delete and reconstruct the created sketch to adapt to the evolving data stream. This brings a huge efficiency cost and temporal information loss. 
Another drawback of \MS is that it needs a predefined number of clusters ($K$) that is hard to determine for clustering of real-world 
 evolving data streams.
% In order to guarantee its clustering performance for adapting to the evolving data stream, the created sketch will be frequently deleted and reconstructed, which may bring in a huge efficiency cost and temporal information loss. 
% Additionally, since the number of final clusters extracted from the sketch needs to be predefined, many representative clusters may not be kept during processing. As a result, its clustering ability will also be constrained under the evolving data stream. 

% \label{subsec:DPT}
% \DPT~\cite{EDMStream:17} is one of the most recent summarizing data structure specifically designed for handling \cef. It is constructed based on the cluster density and can handle all types of \cef as mentioned in section~\ref{sec:background}. Similar to \CFT, \DPT designs a special tree node to represent the temporal clusters, which is called \emph{cluster cell}. Although cluster cell is not as flexible as CF, it can not only be used to locate the corresponding cluster, but also stores its temporal density for building the whole hierarchical structure, which may bring in a better clustering accuracy due to using more comprehensive information.

% \subcompact
% \subsection{Micro Clusters (\MCs)}
% \label{subsec:MCs}
% A variety of \dsc algorithms use \emph{micro clusters} (\MCs) as summarizing data structures, such as \dbstream~\cite{DBStream:16}, \denstream~\cite{DenStream:06} and \clustream~\cite{Clustream:03}. 
% % The number of MC in a \MCs is specified by an input parameter $q$.
% Each \MCs contains $q$ (input parameter) MC as a list, where each MC is an extension of the CF. 

% Figure~\ref{fig:data_structure}(b) illustrates the structure of \MCs. 
% Each MC contains all three key elements (N, LS, SS) in CF, and thus can be also computed incrementally and effectively, especially under the occurrence of concept evolution. The additional components are the sum of the timestamps (LST) and the sum of the squares of the timestamps (SST). 
% Furthermore, each MC maintains a maximum boundary computed related to the mean distance of cluster objects' centroids~\cite{Survey13}.
% %constrains

% % \modify{, and each MC has a maximum
% % boundary computed related to the mean distance of cluster objects' centroids.}
% %how it works
% For each data insertion, an MC with the closest Euclidean distance is selected for potential absorption. Whether to absorb depends on 
% 1) the aforementioned maximum boundary of MC 
% and 2) the distance between the inserted data and the centroid of MC.
% A successful absorption will be followed by updating the summary statistics, or a new MC will be created otherwise.
% %clean
% The oldest MC is deleted if its timestamp is below a given threshold $\delta$ (input parameter), which is deemed to be an outlier and therefore removed. 
% We will discuss more details of outlier detection mechanisms in Section~\ref{subsec:outlier_detection_mechanism}.

% \subcompact
% \subsection{Coreset Tree (\CoreT)}
% \label{subsec:CT}
% A summarising data structure called \CoreT is proposed in the \skm algorithm~\cite{StreamKM++:12}. 
% %principle
% The key idea of constructing a \CoreT is to achieve a $\frac{1}{2}$ reduction in the number of objects by iterative expansion and summary over child nodes.

% As illustrated in Figure~\ref{fig:data_structure}(c). 
% \CoreT is a binary tree in which each node $i$ contains the following components: 
% a set of data points, $E_i$; 
% a representative point of $E_i$, $q_i$; 
% the number of data points, $N_i$; 
% and the sum of squared distances of the data points to $q_i$, $SSE_i$. 
% $E_i$ only has to be stored in the leaf nodes, because the objects of an inner node are implicitly defined as the union of the objects of its child nodes.

% \CoreT cooperates with a bucket set in \skm~\cite{StreamKM++:12} in an iterative \emph{merge-reduce} fashion. The bucket set has a finite amount of buckets and each bucket has a finite capacity of data objects, as specified in two corresponding input parameters.
% Each input data is firstly stored in a bucket of the bucket set, and the data objects in a full bucket will be moved to another bucket in the bucket set. Each pair of two full buckets will be used to construct a \CoreT for object reduction, and $\frac{1}{2}$ results will be stored in another non-empty bucket.

% in order to speed up the sampling non-uniformly during the coreset construction

% \subcompact
% \subsection{Grids (\Grids)}
% \label{subsec:Grids}
% Some \dsc algorithms such as \dstream, perform data summarization through \Grids~\cite{DStream:2007}. That is, by partitioning the n-dimensional feature space into density grid cells.  
% % Taking \dstream as an example,

% As shown in Figure~\ref{fig:data_structure}(d),
% each grid cell is represented by a tuple <$t_g$, $t_m$, $D$, $label$, $status$>.
% $t_g$ is the last time the grid cell was updated, 
% $t_m$ is the last time the grid cell was removed from the hash table that holds the valid grid cells, $D$ is the grid cell density at its last update, $label$ is the class label of the grid cell, and $status$ indicates whether the grid cell is normal or Sporadic (cells with few points are likely outliers).
% %

% % The grid cell is updated online in receiving new data and conducting data cleaning.
% % Specifically, 
% Each new input data point with $d$ dimensions is mapped to the corresponding grid cell based on the summary statistics of previous input data stored in the grid cell, then the clustering is performed based on the density of grid cells. 
% A grid cell is marked as sparse if it does not receive new objects for a long time and marked as dense otherwise. In addition, sporadic grid cells are periodically removed from the list of valid grid cells.

% the grid density is updated upon the mapping of any new input data record into the grid, and the new grid density can be updated based on the difference between the arrival time of the new data record and the last data record



% Continue to explain how clustering is performed based on the grid density...
% \subcompact
% \subsection{Dependency Tree (\DPT)}
% \label{subsec:DPT}
% The \DPT data structure is proposed by the \edmstream~\cite{EDMStream:17} algorithm. Each node in a \DPT is a \emph{cluster-cell} rather than a single point, which is for saving memory space and computation time.
% % A DP tree where each node is a cluster-cell is used to capture the summarization of data streams. 
 
% As shown in Figure~\ref{fig:data_structure} (e), a cluster-cell $c$ summarizes a group of close points, which can be described by a tuple <$s_c$, $\rho^t_c$, $\delta_c^t$> at time $t$. 
% $s_c$ is the seed point. The cluster-cell $c$ seeded by $s_c$ summarizes a set of points whose distance to $s_c$ is less than the distance to any other seed point and is less than or equal to a predefined radius $r$. $\rho^t_c$ is the summarization of all cluster-cell points' timely density at time $t$. $\delta_c^t$ is the dependent distance from $s_c$ to its nearest cluster-cell seed point with higher cluster-cell density. 
% Cluster-cells with relatively low density will not be stored in the DP tree. Instead, they are considered \emph{temporal outliers}.
% % and cached in an \emph{outlier reservoir} (Section~\ref{subsec:outlier_detec_timer}).

% Each new data point may either enter an existing data cell or lead to the creation of a new data cell to contain it. This process is distance-based and can be controlled by $r$ and $s_c$. After data insertion, the \DPT is responsible for manipulating the cluster-cells based on $t$, e.g., dependency update, emergence and decay.
% % and the time-decay value of data is therefore enabled.

% \subcompact
% \subsection{Meyerson Sketch (\MS)}
% \label{subsec:meyerson_sketch}
% The \MS data structure is initially proposed to handle online facility location problems and further extended to clustering data streams by \slkmeans~\cite{SL-KMeans:20}, which is one of the most up-to-date \dsc algorithms. 

% \MS is made up of two main elements, as shown in Figure~\ref{fig:data_structure} (f). 
% The first element is the cluster set, which stores the whole bunch of data inside the cluster and provides convenience for the algorithm to cluster the incoming data stream. However, since the data stream is infinite, it is unrealistic to store every data point in this set. Therefore, \MS gives every cluster set a unique weight decaying with time. Expired clusters are regularly deleted from the set. 
% The second element is the corresponding centre towards the cluster, which is continuously updated according to the cluster set.

% Every time when a new data point comes, \MS will use the cluster centre to calculate the cost of the whole cluster to determine whether it needs to be added to the temporal cluster. If so, it will be added to the corresponding cluster set.

\textbf{Discussion.}
\label{subsec:data_structure_discussion}
In general, hierarchical summarizing data structure improves the clustering efficiency while the partitional structure is expected to bring in better accuracy. 
The unique design of each specific structure, however, leads to further differences in their clustering behaviour.
% makes them behave differently.
% For every specific structure, their unique design and stored information 
% also influence the clustering behavior. 
For example, although \CFT, \CoreT, \DPT and \MCs promise high clustering accuracy due to the comprehensive information stored in their structure, their efficiency behaviour is quite different. Only \CFT and \DPT can guarantee high efficiency due to their flexible updating. \CoreT and \MS are slow due to frequently rebuilding the whole structure while \MCs is slow since it needs to scan every temporal cluster for data insertion. For \Grids, although it may not be as accurate as the other structures, it largely saves its computation time due to eliminating the distance calculation operation. 
% \MS requires maintaining a great number of raw data for clustering, it is hence expected to be the slowest summarizing data structure among others. 

% These summarizing data structures behave quite differently from each other in handling continuous data insertion.
% % The six summarizing data structures discussed above are all applied by \dsc algorithms to summarize and store the cluster information. During the clustering procedure, they will keep updating upon receiving a new data and utilize the stored information to insert the data into a suitable cluster. However, due to the difference in either catalog or main structure, some of them may behave quite differently from others during this common procedure. 
% {First}, 
% for hierarchical-based summarizing data structures such as \CFT, \CFT and \DPT, they do not need to check every existing cluster during data insertion.
% % clustering like those with partition-based catalog. % between parent and children nodes stored in the structure can 
% The hierarchical relationship helps them achieve a more efficient search for a suitable cluster to insert the new data point. 
% {Second},
% unlike many other summarizing data structures, \Grids does not require frequently computing the centre of clusters during data insertion. Instead, a new data point is quickly mapped to the corresponding grid cell based on the computed boundary of the cell. However, the mapping may be inaccurate or outdated overtime, leading to lower clustering accuracy. 
% % Therefore, \Grids spends much less time for data insertion. 
% % However, it leads to inaccurate data insertion and 
% % also due to the usage of boundary rather than distance, 
% % the data insertion may not be so much precise and thus \Grids may not achieve a very great clustering accuracy. 
% {Third}, 
% \MS is the only summarizing data structure that stores every data point of the cluster. Although it tries to constrain the cluster size, more time is wasted on frequently updating the centre of clusters. Moreover, due to a fixed number of clusters, the information of many historical clusters may get erased from the structure, bringing down the clustering accuracy.
% % during the clustering procedure, when more and more data is added into the stored cluster, more time will be wasted on frequently updating the center of clusters, which may result in a low efficiency. Moreover, due to the fix number of clusters, the information of many historical clusters may get erased from the structure, and this may bring down the final clustering accuracy.
% % Grid structure can ensure high efficiency but brings down the accuracy. MS may bring in a terrible clustering efficiency. Hierarchical-based summarizing data structure}